% Options for packages loaded elsewhere
\PassOptionsToPackage{unicode}{hyperref}
\PassOptionsToPackage{hyphens}{url}
%
\documentclass[
]{article}
\usepackage{amsmath,amssymb}
\usepackage{iftex}
\ifPDFTeX
  \usepackage[T1]{fontenc}
  \usepackage[utf8]{inputenc}
  \usepackage{textcomp} % provide euro and other symbols
\else % if luatex or xetex
  \usepackage{unicode-math} % this also loads fontspec
  \defaultfontfeatures{Scale=MatchLowercase}
  \defaultfontfeatures[\rmfamily]{Ligatures=TeX,Scale=1}
\fi
\usepackage{lmodern}
\ifPDFTeX\else
  % xetex/luatex font selection
  \setmainfont[]{DejaVu Serif}
  \setmonofont[]{DejaVu Sans Mono}
\fi
% Use upquote if available, for straight quotes in verbatim environments
\IfFileExists{upquote.sty}{\usepackage{upquote}}{}
\IfFileExists{microtype.sty}{% use microtype if available
  \usepackage[]{microtype}
  \UseMicrotypeSet[protrusion]{basicmath} % disable protrusion for tt fonts
}{}
\makeatletter
\@ifundefined{KOMAClassName}{% if non-KOMA class
  \IfFileExists{parskip.sty}{%
    \usepackage{parskip}
  }{% else
    \setlength{\parindent}{0pt}
    \setlength{\parskip}{6pt plus 2pt minus 1pt}}
}{% if KOMA class
  \KOMAoptions{parskip=half}}
\makeatother
\usepackage{xcolor}
\usepackage[margin=0.85in]{geometry}
\usepackage{color}
\usepackage{fancyvrb}
\newcommand{\VerbBar}{|}
\newcommand{\VERB}{\Verb[commandchars=\\\{\}]}
\DefineVerbatimEnvironment{Highlighting}{Verbatim}{commandchars=\\\{\}}
% Add ',fontsize=\small' for more characters per line
\newenvironment{Shaded}{}{}
\newcommand{\AlertTok}[1]{\textcolor[rgb]{1.00,0.00,0.00}{\textbf{#1}}}
\newcommand{\AnnotationTok}[1]{\textcolor[rgb]{0.38,0.63,0.69}{\textbf{\textit{#1}}}}
\newcommand{\AttributeTok}[1]{\textcolor[rgb]{0.49,0.56,0.16}{#1}}
\newcommand{\BaseNTok}[1]{\textcolor[rgb]{0.25,0.63,0.44}{#1}}
\newcommand{\BuiltInTok}[1]{\textcolor[rgb]{0.00,0.50,0.00}{#1}}
\newcommand{\CharTok}[1]{\textcolor[rgb]{0.25,0.44,0.63}{#1}}
\newcommand{\CommentTok}[1]{\textcolor[rgb]{0.38,0.63,0.69}{\textit{#1}}}
\newcommand{\CommentVarTok}[1]{\textcolor[rgb]{0.38,0.63,0.69}{\textbf{\textit{#1}}}}
\newcommand{\ConstantTok}[1]{\textcolor[rgb]{0.53,0.00,0.00}{#1}}
\newcommand{\ControlFlowTok}[1]{\textcolor[rgb]{0.00,0.44,0.13}{\textbf{#1}}}
\newcommand{\DataTypeTok}[1]{\textcolor[rgb]{0.56,0.13,0.00}{#1}}
\newcommand{\DecValTok}[1]{\textcolor[rgb]{0.25,0.63,0.44}{#1}}
\newcommand{\DocumentationTok}[1]{\textcolor[rgb]{0.73,0.13,0.13}{\textit{#1}}}
\newcommand{\ErrorTok}[1]{\textcolor[rgb]{1.00,0.00,0.00}{\textbf{#1}}}
\newcommand{\ExtensionTok}[1]{#1}
\newcommand{\FloatTok}[1]{\textcolor[rgb]{0.25,0.63,0.44}{#1}}
\newcommand{\FunctionTok}[1]{\textcolor[rgb]{0.02,0.16,0.49}{#1}}
\newcommand{\ImportTok}[1]{\textcolor[rgb]{0.00,0.50,0.00}{\textbf{#1}}}
\newcommand{\InformationTok}[1]{\textcolor[rgb]{0.38,0.63,0.69}{\textbf{\textit{#1}}}}
\newcommand{\KeywordTok}[1]{\textcolor[rgb]{0.00,0.44,0.13}{\textbf{#1}}}
\newcommand{\NormalTok}[1]{#1}
\newcommand{\OperatorTok}[1]{\textcolor[rgb]{0.40,0.40,0.40}{#1}}
\newcommand{\OtherTok}[1]{\textcolor[rgb]{0.00,0.44,0.13}{#1}}
\newcommand{\PreprocessorTok}[1]{\textcolor[rgb]{0.74,0.48,0.00}{#1}}
\newcommand{\RegionMarkerTok}[1]{#1}
\newcommand{\SpecialCharTok}[1]{\textcolor[rgb]{0.25,0.44,0.63}{#1}}
\newcommand{\SpecialStringTok}[1]{\textcolor[rgb]{0.73,0.40,0.53}{#1}}
\newcommand{\StringTok}[1]{\textcolor[rgb]{0.25,0.44,0.63}{#1}}
\newcommand{\VariableTok}[1]{\textcolor[rgb]{0.10,0.09,0.49}{#1}}
\newcommand{\VerbatimStringTok}[1]{\textcolor[rgb]{0.25,0.44,0.63}{#1}}
\newcommand{\WarningTok}[1]{\textcolor[rgb]{0.38,0.63,0.69}{\textbf{\textit{#1}}}}
\usepackage{longtable,booktabs,array}
\usepackage{calc} % for calculating minipage widths
% Correct order of tables after \paragraph or \subparagraph
\usepackage{etoolbox}
\makeatletter
\patchcmd\longtable{\par}{\if@noskipsec\mbox{}\fi\par}{}{}
\makeatother
% Allow footnotes in longtable head/foot
\IfFileExists{footnotehyper.sty}{\usepackage{footnotehyper}}{\usepackage{footnote}}
\makesavenoteenv{longtable}
\setlength{\emergencystretch}{3em} % prevent overfull lines
\providecommand{\tightlist}{%
  \setlength{\itemsep}{0pt}\setlength{\parskip}{0pt}}
\setcounter{secnumdepth}{-\maxdimen} % remove section numbering
\ifLuaTeX
  \usepackage{selnolig}  % disable illegal ligatures
\fi
\usepackage{bookmark}
\IfFileExists{xurl.sty}{\usepackage{xurl}}{} % add URL line breaks if available
\urlstyle{same}
\hypersetup{
  hidelinks,
  pdfcreator={LaTeX via pandoc}}

\author{}
\date{}

\begin{document}

{
\setcounter{tocdepth}{3}
\tableofcontents
}
\section{Hardware-Aware TLS Identity Binding
Profile}\label{hardware-aware-tls-identity-binding-profile}

Draft v0.1.1

\subsection{Abstract}\label{abstract}

This document defines a repository profile for binding upper-layer
identity and authorization material to an accepted TLS 1.3 session and
to post-handshake platform-attestation facts. The profile is
application-protocol neutral at L0 through L2. Agent2Agent / AGTP is the
first reference target for the upper layers, not a replacement for the
trust model described here.

The core rule is simple: a verifier accepts a peer only when the peer's
authenticated identity, session binding, freshness state, and local
policy all refer to the same intended session, platform, service, agent,
task, and authority boundary.

\subsection{1. Status}\label{status}

This is a draft security-hardening profile for implementation review. It
is not an IETF consensus document and it does not define a new TLS
handshake, a new attestation evidence format, a new identity provider,
or a new application core protocol.

Within this repository, this Markdown file is the normative source for
the profile. Rendered specifications, notes, reports, and tests are
derived from it unless they explicitly say otherwise.

The key words ``MUST'', ``MUST NOT'', ``REQUIRED'', ``SHALL'', ``SHALL
NOT'', ``SHOULD'', ``SHOULD NOT'', ``RECOMMENDED'', ``NOT RECOMMENDED'',
``MAY'', and ``OPTIONAL'' are to be interpreted as described in BCP 14,
RFC 2119 and RFC 8174 when, and only when, they appear in all capitals.

\subsection{2. Terms}\label{terms}

Hardware-aware TLS means ordinary TLS 1.3 plus post-handshake platform
attestation and session binding. Platform attestation does not replace
the TLS handshake and does not authenticate the platform before a TLS
channel exists.

The older shorthand \texttt{aTLS} is reserved for existing Cocos
implementation code paths, package names, or historical text. New
profile text should use ``hardware-aware TLS'' unless it refers to that
legacy surface.

Identity Grant means an authority statement issued by a Manager or
another locally trusted policy authority. It authorizes upper-layer
identity, task, and authorization values.

Session Binding Statement means a holder-of-key proof that binds one
verified Identity Grant to one accepted TLS or exported-authenticator
session. It does not authorize service, tenant, task, scope, resource,
or capability values.

Expected values are verifier-local policy inputs. Observed values are
values extracted from authenticated grants, session-bound statements,
attestation claims, trusted manifests, or locally derived request state.
Observed values do not become expected values without a trusted local
policy decision.

\subsection{3. Scope}\label{scope}

This profile covers the checks needed before an application treats an
accepted TLS peer as the intended platform, service, agent, task, or
authorized actor. It focuses on fail-closed binding across three classes
of facts:

\begin{enumerate}
\def\labelenumi{\arabic{enumi}.}
\tightlist
\item
  accepted TLS and post-handshake attestation facts;
\item
  authenticated identity and authorization material;
\item
  verifier-local expected policy.
\end{enumerate}

Out of scope:

\begin{itemize}
\tightlist
\item
  selecting a specific identity provider;
\item
  standardizing a concrete wire-token format;
\item
  changing attestation evidence formats;
\item
  changing lower-layer TLS or post-handshake attestation wire protocols;
\item
  proving a complete end-to-end authorization model;
\item
  defining generic OAuth, OIDC, JWT, JWS, CWT, or COSE behavior except
  where this profile relies on them for fail-closed acceptance.
\end{itemize}

A requirement belongs in this profile when omitting it could cause any
of these acceptance failures:

\begin{itemize}
\tightlist
\item
  a binding from the wrong live TLS or attestation session is accepted;
\item
  the right platform is accepted for the wrong service, tenant, agent,
  task, or authority boundary;
\item
  replay, stale state, missing state, or required-mode downgrade is
  accepted;
\item
  peer-selected metadata becomes verifier policy;
\item
  a security result, authorization decision, or verification result is
  reused outside the session or request for which it was produced.
\end{itemize}

\subsection{4. Relationship to existing
specifications}\label{relationship-to-existing-specifications}

Most building blocks are already standardized. This profile composes
them and adds repository-specific fail-closed requirements for identity,
session binding, replay, canonical references, local policy comparison,
and cache safety.

\begin{longtable}[]{@{}
  >{\raggedright\arraybackslash}p{(\columnwidth - 4\tabcolsep) * \real{0.3333}}
  >{\raggedright\arraybackslash}p{(\columnwidth - 4\tabcolsep) * \real{0.3333}}
  >{\raggedright\arraybackslash}p{(\columnwidth - 4\tabcolsep) * \real{0.3333}}@{}}
\toprule\noalign{}
\begin{minipage}[b]{\linewidth}\raggedright
Area
\end{minipage} & \begin{minipage}[b]{\linewidth}\raggedright
Existing specification
\end{minipage} & \begin{minipage}[b]{\linewidth}\raggedright
Profile decision
\end{minipage} \\
\midrule\noalign{}
\endhead
\bottomrule\noalign{}
\endlastfoot
Normative keywords & BCP 14: RFC 2119 and RFC 8174 & Keywords apply to
this repository profile only. \\
OAuth roles and authorization vocabulary & RFC 6749 and RFC 9396 &
Manager, Agent, relying party, \texttt{scope}, \texttt{resource}, and
\texttt{authorization\_details} are mapped by this profile. \\
JWT/JWS syntax and safety & RFC 7519, RFC 7515, and RFC 8725 & The
two-token Identity Grant plus Session Binding Statement model,
\texttt{agtp\_type}, \texttt{agtp\_version}, \texttt{grant\_hash}, and
the mandatory claim set are profile choices. \\
Proof-of-possession pattern & RFC 7800 and RFC 8705 & The Agent
confirmation-key rule follows the PoP pattern. The exact Session Binding
Statement claims are profile-specific. \\
TLS channel binding and exported authenticators & RFC 9266 and RFC 9261
& The profile consumes accepted lower-layer binding facts. It does not
define a new TLS exporter or exported-authenticator format. \\
Remote attestation roles & RFC 9334 & The attestation-binder hash,
accepted-evidence policy, and L1/L2 wiring are profile details. \\
CWT/COSE & RFC 8392 and RFC 9052 & CWT/COSE is an alternative encoding
for the same semantics, not a different trust model. \\
HTTP response caching & RFC 9111 & Response-cache vocabulary is reused.
Security-state objects and verification results are not cacheable
acceptance evidence. \\
OIDC & OpenID Connect Core & OIDC is vocabulary only. This profile does
not require OIDC. \\
L0-L6, Identity Grant, Session Binding Statement, canonical references &
Not standardized as one combined profile & These are local names for
reviewable fail-closed identity binding. \\
\end{longtable}

\subsection{5. Threat model}\label{threat-model}

The attacker can observe, replay, reorder, or relay network messages.
The attacker can provide malicious application metadata. The attacker
can run another Agent on the same machine or in the same deployment
environment.

In the Alice-side example, malware may influence peer-provided metadata
or try to reuse old grants or binding statements. It does not fully
control Alice's verifier process, local policy store, OS sandbox,
keychain, or trusted Manager key configuration.

The lower layers are assumed to implement TLS 1.3 validation,
exported-authenticator validation, and platform-evidence appraisal
correctly. This profile defines the identity and policy material that
must be bound to those lower-layer facts before the application accepts
the peer as the intended Agent.

Out of scope for this threat model:

\begin{itemize}
\tightlist
\item
  a fully compromised verifier process or local policy store;
\item
  a malicious or compromised trusted Manager or policy authority;
\item
  compromise of all local secret storage, OS isolation, or keychain
  controls;
\item
  denial of service;
\item
  side-channel leakage outside the identity-binding path.
\end{itemize}

\begin{longtable}[]{@{}
  >{\raggedright\arraybackslash}p{(\columnwidth - 2\tabcolsep) * \real{0.5000}}
  >{\raggedright\arraybackslash}p{(\columnwidth - 2\tabcolsep) * \real{0.5000}}@{}}
\toprule\noalign{}
\begin{minipage}[b]{\linewidth}\raggedright
Threat
\end{minipage} & \begin{minipage}[b]{\linewidth}\raggedright
Required design consequence
\end{minipage} \\
\midrule\noalign{}
\endhead
\bottomrule\noalign{}
\endlastfoot
Network relay or borrowed attestation & Bind the verified grant to the
accepted TLS session, including endpoint-key hash, request-context hash,
and attestation-binder hash when present. \\
Replay of old grants or bindings & Use \texttt{iat}, \texttt{exp},
unique IDs, and one-shot nonce or binding state. Required-mode
deployments fail closed when replay state is unavailable. \\
Same-machine wrong-Agent & Keep Manager keys and Agent confirmation keys
in separate trust domains. Accept only a Session Binding Statement
signed by the grant-authorized confirmation key or by a key explicitly
authorized by local policy. \\
Peer-provided metadata injection & Build accepted identity only from
authenticated grants, session-bound statements, and local expected
policy. \\
Service, tenant, deployment, task, or capability diversion & Compare L3
through L6 values against local expected values. Required layers fail
closed when expected values are missing or ambiguous. \\
Token substitution between grant and binding & Carry a domain-separated
\texttt{grant\_hash} over the exact signed grant bytes. Reject
mismatched grant and binding pairs. \\
Key rotation, stale key, or revocation failure & Reject unknown, stale,
retired, or revoked Manager keys, Agent binding keys, and grant IDs.
Remote key sources need freshness and failure behavior. \\
Cache confusion or cross-authority leakage & Do not reuse identity
tokens, binding statements, attestation evidence, authorization
decisions, or verification results as acceptance evidence. \\
Gateway route confusion & In gateway-terminated deployments,
authenticate the gateway-to-Agent route. Gateway session binding alone
proves only the gateway endpoint. \\
Semantic alias or fuzzy-match confusion & Use canonical identifiers for
decision-sensitive semantic values. Fuzzy matching, alias repair, or
model interpretation is not an acceptance authority. \\
\end{longtable}

If Alice's verifier, local expected-policy source, or trusted key store
is fully compromised, token design alone cannot recover security. Those
cases require host isolation, policy-file integrity, secret isolation,
process isolation, and operational controls outside this profile.

\subsection{6. Layer model}\label{layer-model}

The layer model separates lower-layer session and attestation binding
from upper-layer deployment, agent, task, and authorization policy.
Other profiles may collapse, rename, or extend these layers if they
preserve the same fail-closed checks.

\begin{longtable}[]{@{}
  >{\raggedright\arraybackslash}p{(\columnwidth - 4\tabcolsep) * \real{0.3333}}
  >{\raggedright\arraybackslash}p{(\columnwidth - 4\tabcolsep) * \real{0.3333}}
  >{\raggedright\arraybackslash}p{(\columnwidth - 4\tabcolsep) * \real{0.3333}}@{}}
\toprule\noalign{}
\begin{minipage}[b]{\linewidth}\raggedright
Layer
\end{minipage} & \begin{minipage}[b]{\linewidth}\raggedright
Verification target
\end{minipage} & \begin{minipage}[b]{\linewidth}\raggedright
Main failure class
\end{minipage} \\
\midrule\noalign{}
\endhead
\bottomrule\noalign{}
\endlastfoot
L0 & Live TLS channel & MITM or session confusion \\
L1 & Attested platform validity & Fake, malformed, or untrusted platform
evidence \\
L2 & Attestation-to-channel or authenticator-to-channel binding & Relay,
replay, or borrowed evidence \\
L3 & Intended service, tenant, deployment, or environment & Service or
tenant diversion \\
L4 & Intended workload, process, or agent & Same-machine wrong-Agent \\
L5 & Intended task, thread, context, or delegation & Cross-task or
context diversion \\
L6 & Intended authorization or capability policy & Confused deputy or
privilege escalation \\
\end{longtable}

L0 through L2 can be tested with transport, attestation, and
implementation regressions. L3 through L6 require explicit policy inputs
before a verifier or application layer can enforce them consistently.

\subsection{7. Application-protocol
boundary}\label{application-protocol-boundary}

TLS and the hardware-aware profile are responsible for L0 through L2:

\begin{itemize}
\tightlist
\item
  authenticating the live TLS channel;
\item
  appraising attested platform or VM evidence;
\item
  binding authenticator or attestation material to the accepted TLS
  session.
\end{itemize}

An application profile carries or references the upper-layer identity
and policy material needed for L3 through L6:

\begin{itemize}
\tightlist
\item
  intended service, tenant, deployment, or environment;
\item
  intended workload, process, or Agent;
\item
  intended task, thread, context, or delegation;
\item
  authorization or capability policy.
\end{itemize}

Relay defense remains an L2 session-binding question. Diversion and
wrong-Agent prevention require L3 and L4 policy comparison. Task and
capability checks then continue at L5 and L6.

Current application-protocol drafts do not by themselves provide all
semantic binding needed for service, tenant, Agent, task, capability,
canonical-reference, replay, and local-policy checks. The application
protocol may carry authenticated policy inputs, such as an Identity
Grant, and bind them to the accepted TLS session through a Session
Binding Statement. The verifier still compares the resulting observed
values with local expected values before accepting the peer.

\begin{longtable}[]{@{}
  >{\raggedright\arraybackslash}p{(\columnwidth - 6\tabcolsep) * \real{0.2500}}
  >{\raggedright\arraybackslash}p{(\columnwidth - 6\tabcolsep) * \real{0.2500}}
  >{\raggedright\arraybackslash}p{(\columnwidth - 6\tabcolsep) * \real{0.2500}}
  >{\raggedright\arraybackslash}p{(\columnwidth - 6\tabcolsep) * \real{0.2500}}@{}}
\toprule\noalign{}
\begin{minipage}[b]{\linewidth}\raggedright
Review item
\end{minipage} & \begin{minipage}[b]{\linewidth}\raggedright
Question
\end{minipage} & \begin{minipage}[b]{\linewidth}\raggedright
Profile requirement
\end{minipage} & \begin{minipage}[b]{\linewidth}\raggedright
Negative-test anchor
\end{minipage} \\
\midrule\noalign{}
\endhead
\bottomrule\noalign{}
\endlastfoot
SP-01 & Is identity material bound to the accepted TLS and attestation
session? & Verify the Session Binding Statement against endpoint key,
request/exporter context, attestation binder when present, and replay
state. & \texttt{hwtls-id-profile-relay-001}; attestation-binder
mismatch \\
SP-02 & Which values are local policy and which are peer claims? &
Compare service, tenant, deployment, agent, task, scope, resource, and
authorization values against local expected policy. &
\texttt{hwtls-id-profile-diversion-001};
\texttt{hwtls-id-profile-wrong-agent-001};
\texttt{hwtls-id-profile-binding-confusion-001} \\
SP-03 & Which freshness values are one-shot? & Reject repeated Session
Binding Statement IDs, nonces, or task-binding values when one-shot use
is required. & \texttt{hwtls-id-profile-replay-001} \\
SP-04 & What happens when grant or binding material is missing,
unsupported, substituted, or partly verified? & Required mode fails
closed. Neither Identity Grant nor Session Binding Statement is
sufficient alone. & \texttt{hwtls-id-profile-downgrade-001};
nested-token substitution \\
SP-05 & Which responses are safe for shared caching? & Treat
caller-dependent responses as \texttt{private} with adequate
partitioning or \texttt{no-store}. Do not cache security inputs or
verification results as acceptance evidence. & caller-dependent cache
leak \\
SP-06 & Are semantic references canonical? & Reject non-canonical,
ambiguous, fuzzy-matched, or receiver-repaired decision values. &
semantic alias confusion \\
SP-07 & Are Manager and Agent keys separated? & Reject Manager keys used
as Agent confirmation keys and Agent keys used as Manager signing keys.
& key-role confusion \\
SP-08 & Is gateway routing authenticated? & In gateway mode, bind the
gateway endpoint and the gateway-to-Agent route separately. & gateway
route confusion \\
\end{longtable}

\subsection{8. Authority and key
separation}\label{authority-and-key-separation}

One implementation profile used in this repository is:

\begin{Shaded}
\begin{Highlighting}[]
\NormalTok{Identity Grant + hardware{-}aware TLS + OAuth/OIDC{-}style semantics + JWT/JWS encoding}
\end{Highlighting}
\end{Shaded}

OAuth and OIDC provide claim semantics and review vocabulary. JWT/JWS
provides one signed-token encoding. CWT/COSE provides a compact binary
encoding for the same trust model.

The profile keeps three key roles separate:

\begin{longtable}[]{@{}
  >{\raggedright\arraybackslash}p{(\columnwidth - 2\tabcolsep) * \real{0.5000}}
  >{\raggedright\arraybackslash}p{(\columnwidth - 2\tabcolsep) * \real{0.5000}}@{}}
\toprule\noalign{}
\begin{minipage}[b]{\linewidth}\raggedright
Key role
\end{minipage} & \begin{minipage}[b]{\linewidth}\raggedright
Purpose
\end{minipage} \\
\midrule\noalign{}
\endhead
\bottomrule\noalign{}
\endlastfoot
TLS endpoint key & Proves possession for the accepted TLS or
exported-authenticator endpoint. \\
Agent binding key & Signs the Session Binding Statement that binds a
verified grant to the accepted TLS session. \\
Manager or policy-authority key & Signs Identity Grants that authorize
intended deployment, Agent, task, or capability values. \\
\end{longtable}

These keys may be related by deployment policy, but they MUST NOT be
treated as the same key by default. The Manager key is a token-signing
or policy-authority key, not a TLS endpoint key.

Manager signing keys MUST NOT be accepted as Agent confirmation keys.
Agent confirmation keys MUST NOT be accepted as Manager or
policy-authority signing keys. A deployment MAY explicitly authorize an
endpoint key for session binding, but that authorization must come from
the verified grant or from local policy; it must not be inferred from
the TLS session alone.

The JWS protected header \texttt{kid} is only a key-selection hint.
Acceptance also requires issuer, audience, algorithm allow-list, key
status, key use, profile version, token type, time validity, and local
policy checks. The signing-method allow-list MUST NOT include
\texttt{none}.

\subsection{9. Identity Grant}\label{identity-grant}

An Identity Grant authorizes the intended upper-layer subject. The
initial wire profile uses JWT/JWS and OAuth/OIDC-style claim names where
they fit. Other encodings, including CWT/COSE or a signed manifest, are
acceptable only if they preserve the same authority, freshness,
canonicalization, and session-binding rules.

An Identity Grant contains the deployment-specific equivalent of:

\begin{itemize}
\tightlist
\item
  profile token type, \texttt{agtp\_type=agtp.identity-grant} in the
  current JWT/JWS implementation;
\item
  profile version, \texttt{agtp\_version=1} in the current JWT/JWS
  implementation;
\item
  issuer, \texttt{iss};
\item
  subject, \texttt{sub};
\item
  audience, \texttt{aud};
\item
  unique grant ID, \texttt{jti} for JWT/JWS or \texttt{cti} for
  CWT/COSE;
\item
  issued-at time, \texttt{iat};
\item
  expiration time, \texttt{exp};
\item
  Agent confirmation key, \texttt{cnf.kid} in the initial JWT/JWS
  profile;
\item
  optional explicitly authorized endpoint keys;
\item
  service, tenant, deployment, or environment values when L3 is
  required;
\item
  workload, process, or Agent identity when L4 is required;
\item
  computation, task, thread, or delegation IDs when L5 is required;
\item
  canonical intent, capability, or ontology references when those values
  are decision-sensitive;
\item
  scopes, resources, or authorization details when L6 is required;
\item
  issuer key ID or key version when needed for key rotation.
\end{itemize}

The Agent is not the authority for service, tenant, deployment, task,
scope, resource, or capability values. Those values are accepted only
when they appear in a grant verified under a locally trusted Manager or
policy-authority key and then pass local policy comparison.

\subsection{10. Session Binding
Statement}\label{session-binding-statement}

A Session Binding Statement proves that the holder of the confirmation
key named in the verified grant bound that exact grant to the accepted
TLS or exported-authenticator session. It does not authorize identity or
capability values.

A generic accepted endpoint key is not sufficient unless the verified
grant or local attestation policy explicitly binds that key to the same
Agent identity.

A Session Binding Statement contains the deployment-specific equivalent
of:

\begin{itemize}
\tightlist
\item
  profile token type, \texttt{agtp\_type=agtp.session-binding} in the
  current JWT/JWS implementation;
\item
  profile version, \texttt{agtp\_version=1} in the current JWT/JWS
  implementation;
\item
  unique binding statement ID, \texttt{jti};
\item
  audience or relying-service ID, \texttt{aud};
\item
  issued-at time, \texttt{iat};
\item
  expiration time, \texttt{exp};
\item
  \texttt{grant\_hash};
\item
  \texttt{leaf\_public\_key\_sha256};
\item
  \texttt{request\_context\_sha256};
\item
  \texttt{attestation\_binder\_sha256} when accepted
  attestation-to-channel evidence is present;
\item
  one-shot nonce or unique binding value.
\end{itemize}

Binding means receiver-verifiable linkage to the accepted session and
accepted evidence. The statement does not need to embed a TLS session
identifier when the verifier can deterministically compare endpoint-key,
exporter or request context, attestation-binder, nonce, and replay
fields.

The grant hash is computed over an unambiguous byte string. For the
initial encodings:

\begin{Shaded}
\begin{Highlighting}[]
\NormalTok{SHA{-}256("agtp.identity{-}grant.jwt.v1" || NUL || exact{-}signed{-}grant{-}bytes)}
\NormalTok{SHA{-}256("agtp.identity{-}grant.cwt.v1" || NUL || exact{-}signed{-}grant{-}bytes)}
\end{Highlighting}
\end{Shaded}

For JSON-based formats, the verifier SHOULD hash the exact signed bytes.
It MUST NOT compute \texttt{grant\_hash} by re-serializing JSON in the
acceptance path. Canonical CBOR/COSE is acceptable when the encoding
profile defines the canonical form.

When accepted attestation evidence includes an attestation-to-channel
binder, \texttt{attestation\_binder\_sha256} is REQUIRED in the Session
Binding Statement and must match the accepted binder. When no such
binder exists in the lower-layer session, the field is absent unless
deployment policy requires it.

\subsection{11. Verification model}\label{verification-model}

The verifier performs authority checks, session-binding checks, replay
checks, and policy comparison in that order. The final accepted
assertion is built only from verified material.

\begin{enumerate}
\def\labelenumi{\arabic{enumi}.}
\tightlist
\item
  Verify the Identity Grant under a trusted Manager or policy-authority
  key.
\item
  Check grant issuer, audience, profile version, token type, algorithm,
  key status, expiration, issued-at time, grant ID, and required scope
  or resource fields.
\item
  Verify the Session Binding Statement signature.
\item
  Check that the statement signer is the grant confirmation key or
  another endpoint key explicitly authorized by the verified grant or
  local policy.
\item
  Verify that \texttt{grant\_hash} matches the exact verified Identity
  Grant bytes.
\item
  Verify that the statement audience matches the relying service or
  client.
\item
  Reject arbitrary accepted endpoint keys that are not explicitly bound
  to the same Agent identity by the verified grant or local attestation
  policy.
\item
  Compare the statement binding fields with the accepted TLS or
  exported-authenticator session.
\item
  Enforce replay policy for grant IDs, binding IDs, task-binding values,
  and nonces when one-shot use is required.
\item
  Construct the internal assertion only from the verified grant and
  binding statement.
\item
  Compare the assertion with local expected policy for every required
  layer.
\end{enumerate}

\texttt{identitypolicy.ValidateAssertion} remains a comparator and
binding freshness checker. It does not parse wire tokens, verify
signatures, select keys, rotate keys, check revocation, or own replay
storage. Those checks belong to the component that authenticates
external identity material before it returns an observed assertion.

\subsection{12. Policy model}\label{policy-model}

Expected values and observed values stay separate.

\begin{itemize}
\tightlist
\item
  Expected values come from local policy, operator configuration, a
  trusted registry, a Manager, computation state, or a policy engine.
\item
  Observed values come from a session-bound assertion built from
  authenticated grants, Session Binding Statements, attestation
  evidence, CoRIM metadata, EAT claims, agent manifests, authenticated
  request metadata, or authorization tokens.
\item
  Verification compares observed values against expected values.
\item
  Required layers fail closed when an expected value is missing,
  ambiguous, or only peer supplied.
\end{itemize}

A minimal policy object can be shaped as follows:

\begin{Shaded}
\begin{Highlighting}[]
\FunctionTok{identity\_policy}\KeywordTok{:}
\AttributeTok{  }\FunctionTok{mode}\KeywordTok{:}\AttributeTok{ }\StringTok{"disabled"}\CommentTok{        \# disabled | required}
\AttributeTok{  }\FunctionTok{set\_mode}\KeywordTok{:}\AttributeTok{ }\StringTok{"contains\_all"}\CommentTok{ \# contains\_all | exact}

\AttributeTok{  }\FunctionTok{require}\KeywordTok{:}
\AttributeTok{    }\FunctionTok{l3}\KeywordTok{:}\AttributeTok{ }\CharTok{false}
\AttributeTok{    }\FunctionTok{l4}\KeywordTok{:}\AttributeTok{ }\CharTok{false}
\AttributeTok{    }\FunctionTok{l5}\KeywordTok{:}\AttributeTok{ }\CharTok{false}
\AttributeTok{    }\FunctionTok{l6}\KeywordTok{:}\AttributeTok{ }\CharTok{false}

\AttributeTok{  }\FunctionTok{expected}\KeywordTok{:}
\AttributeTok{    }\FunctionTok{service}\KeywordTok{:}\AttributeTok{ }\StringTok{""}
\AttributeTok{    }\FunctionTok{tenant}\KeywordTok{:}\AttributeTok{ }\StringTok{""}
\AttributeTok{    }\FunctionTok{deployment}\KeywordTok{:}\AttributeTok{ }\StringTok{""}
\AttributeTok{    }\FunctionTok{environment}\KeywordTok{:}\AttributeTok{ }\StringTok{""}
\AttributeTok{    }\FunctionTok{workload}\KeywordTok{:}\AttributeTok{ }\StringTok{""}
\AttributeTok{    }\FunctionTok{agent}\KeywordTok{:}\AttributeTok{ }\StringTok{""}
\AttributeTok{    }\FunctionTok{agent\_public\_key}\KeywordTok{:}\AttributeTok{ }\StringTok{""}
\AttributeTok{    }\FunctionTok{computation\_id}\KeywordTok{:}\AttributeTok{ }\StringTok{""}
\AttributeTok{    }\FunctionTok{task\_id}\KeywordTok{:}\AttributeTok{ }\StringTok{""}
\AttributeTok{    }\FunctionTok{thread\_id}\KeywordTok{:}\AttributeTok{ }\StringTok{""}
\AttributeTok{    }\FunctionTok{delegation\_id}\KeywordTok{:}\AttributeTok{ }\StringTok{""}
\AttributeTok{    }\FunctionTok{intent\_ref}\KeywordTok{:}\AttributeTok{ }\StringTok{""}
\AttributeTok{    }\FunctionTok{capability\_ref}\KeywordTok{:}\AttributeTok{ }\StringTok{""}
\AttributeTok{    }\FunctionTok{ontology\_id}\KeywordTok{:}\AttributeTok{ }\StringTok{""}
\AttributeTok{    }\FunctionTok{scopes}\KeywordTok{:}\AttributeTok{ }\KeywordTok{[]}
\AttributeTok{    }\FunctionTok{resources}\KeywordTok{:}\AttributeTok{ }\KeywordTok{[]}
\AttributeTok{    }\FunctionTok{authorization\_details}\KeywordTok{:}\AttributeTok{ }\KeywordTok{[]}
\end{Highlighting}
\end{Shaded}

After external identity material has been authenticated, the verifier
can build an internal assertion:

\begin{Shaded}
\begin{Highlighting}[]
\FunctionTok{identity\_assertion}\KeywordTok{:}
\AttributeTok{  }\FunctionTok{issuer}\KeywordTok{:}\AttributeTok{ }\StringTok{"manager{-}or{-}policy{-}engine"}

\AttributeTok{  }\FunctionTok{values}\KeywordTok{:}
\AttributeTok{    }\FunctionTok{service}\KeywordTok{:}\AttributeTok{ }\StringTok{"payment{-}agent"}
\AttributeTok{    }\FunctionTok{tenant}\KeywordTok{:}\AttributeTok{ }\StringTok{"tenant{-}a"}
\AttributeTok{    }\FunctionTok{deployment}\KeywordTok{:}\AttributeTok{ }\StringTok{"prod"}
\AttributeTok{    }\FunctionTok{environment}\KeywordTok{:}\AttributeTok{ }\StringTok{"asia{-}northeast1"}
\AttributeTok{    }\FunctionTok{workload}\KeywordTok{:}\AttributeTok{ }\StringTok{"settlement{-}worker"}
\AttributeTok{    }\FunctionTok{agent}\KeywordTok{:}\AttributeTok{ }\StringTok{"agent{-}a"}
\AttributeTok{    }\FunctionTok{agent\_public\_key}\KeywordTok{:}\AttributeTok{ }\StringTok{"sha256:..."}
\AttributeTok{    }\FunctionTok{computation\_id}\KeywordTok{:}\AttributeTok{ }\StringTok{"cmp{-}..."}
\AttributeTok{    }\FunctionTok{task\_id}\KeywordTok{:}\AttributeTok{ }\StringTok{"task{-}..."}
\AttributeTok{    }\FunctionTok{thread\_id}\KeywordTok{:}\AttributeTok{ }\StringTok{"thread{-}..."}
\AttributeTok{    }\FunctionTok{delegation\_id}\KeywordTok{:}\AttributeTok{ }\StringTok{"delegation{-}..."}
\AttributeTok{    }\FunctionTok{intent\_ref}\KeywordTok{:}\AttributeTok{ }\StringTok{"intent:monthly{-}settlement"}
\AttributeTok{    }\FunctionTok{capability\_ref}\KeywordTok{:}\AttributeTok{ }\StringTok{"capability:orders{-}read"}
\AttributeTok{    }\FunctionTok{ontology\_id}\KeywordTok{:}\AttributeTok{ }\StringTok{"ontology:payments:v1"}
\AttributeTok{    }\FunctionTok{scopes}\KeywordTok{:}
\AttributeTok{      }\KeywordTok{{-}}\AttributeTok{ }\StringTok{"orders:read"}
\AttributeTok{    }\FunctionTok{resources}\KeywordTok{:}
\AttributeTok{      }\KeywordTok{{-}}\AttributeTok{ }\StringTok{"orders"}
\AttributeTok{    }\FunctionTok{authorization\_details}\KeywordTok{:}
\AttributeTok{      }\KeywordTok{{-}}\AttributeTok{ }\StringTok{"purpose:monthly{-}settlement"}

\AttributeTok{  }\FunctionTok{binding}\KeywordTok{:}
\AttributeTok{    }\FunctionTok{leaf\_public\_key\_sha256}\KeywordTok{:}\AttributeTok{ }\StringTok{"..."}
\AttributeTok{    }\FunctionTok{request\_context\_sha256}\KeywordTok{:}\AttributeTok{ }\StringTok{"..."}
\AttributeTok{    }\FunctionTok{attestation\_binder\_sha256}\KeywordTok{:}\AttributeTok{ }\StringTok{"..."}
\AttributeTok{    }\FunctionTok{nonce}\KeywordTok{:}\AttributeTok{ }\StringTok{"..."}
\AttributeTok{    }\FunctionTok{issued\_at}\KeywordTok{:}\AttributeTok{ }\StringTok{"..."}
\AttributeTok{    }\FunctionTok{expires\_at}\KeywordTok{:}\AttributeTok{ }\StringTok{"..."}
\end{Highlighting}
\end{Shaded}

This assertion is not a wire format. It is a local representation used
after the caller has authenticated external identity material.
\texttt{Assertion.Issuer} is informational unless the observed-identity
implementation has already verified the corresponding issuer and trust
anchor.

\subsection{13. L3 through L6 policy
fields}\label{l3-through-l6-policy-fields}

\begin{longtable}[]{@{}
  >{\raggedright\arraybackslash}p{(\columnwidth - 4\tabcolsep) * \real{0.3333}}
  >{\raggedright\arraybackslash}p{(\columnwidth - 4\tabcolsep) * \real{0.3333}}
  >{\raggedright\arraybackslash}p{(\columnwidth - 4\tabcolsep) * \real{0.3333}}@{}}
\toprule\noalign{}
\begin{minipage}[b]{\linewidth}\raggedright
Layer
\end{minipage} & \begin{minipage}[b]{\linewidth}\raggedright
Required question
\end{minipage} & \begin{minipage}[b]{\linewidth}\raggedright
Common expected values
\end{minipage} \\
\midrule\noalign{}
\endhead
\bottomrule\noalign{}
\endlastfoot
L3 & Is this the intended service, tenant, deployment, or environment? &
service, tenant, deployment, environment, region, CoRIM or EAT
deployment claims \\
L4 & Is this the intended workload, process, or Agent? & workload ID,
Agent ID, process or binary hash, config or policy hash, Agent public
key, routing target \\
L5 & Is this the intended task, thread, context, or delegation? &
computation ID, task ID, thread ID, request context, delegation ID,
callback binding, one-shot task state \\
L6 & Is the action authorized for this peer and context? & scope,
resource, authorization details, capability token, policy decision, user
consent, tool policy \\
\end{longtable}

OIDC and OAuth-style names are useful analogies:

\begin{longtable}[]{@{}
  >{\raggedright\arraybackslash}p{(\columnwidth - 2\tabcolsep) * \real{0.5000}}
  >{\raggedright\arraybackslash}p{(\columnwidth - 2\tabcolsep) * \real{0.5000}}@{}}
\toprule\noalign{}
\begin{minipage}[b]{\linewidth}\raggedright
Layer
\end{minipage} & \begin{minipage}[b]{\linewidth}\raggedright
OAuth/OIDC-style analogue
\end{minipage} \\
\midrule\noalign{}
\endhead
\bottomrule\noalign{}
\endlastfoot
L3 & issuer, audience, tenant, hosted domain, deployment claim,
environment claim \\
L4 & subject, client ID, workload identity, actor claim, confirmation
key \\
L5 & nonce, state, transaction ID, request object, delegation token \\
L6 & scope, resource indicator, authorization details, policy decision,
consent \\
\end{longtable}

These analogies are not normative requirements. A deployment can source
the same expected values from configuration, a registry, CoRIM metadata,
an EAT claim, an agent manifest, or an application policy engine.

\subsection{14. Canonical semantic
references}\label{canonical-semantic-references}

Some L5 and L6 inputs are semantic references rather than raw labels:

\begin{itemize}
\tightlist
\item
  \texttt{ontologyId}: stable identifier for the vocabulary or policy
  namespace;
\item
  \texttt{intentRef}: stable identifier for the intended task, intent,
  or action family;
\item
  \texttt{capabilityRef}: stable identifier for an authorization
  capability or resource-action class.
\end{itemize}

These values are profile reference identifiers. They are not
natural-language prompts, display names, or peer-selected aliases.

Decision-sensitive semantic values MUST be canonical before they appear
in an Identity Grant, Session Binding Statement, local expected policy,
or authorization decision. The relying client, Manager, or local policy
authority resolves aliases before issuance or local policy comparison.

Receivers MUST validate canonical identifiers deterministically against
receiver-side policy, task state, trusted registry data, signed profile
data, or the underlying interaction/security protocol.

Receivers MUST NOT normalize peer-provided decision-sensitive values
during the final acceptance path. They MUST NOT turn a peer-provided
alias, display label, natural-language phrase, case variant, URI
variant, or model interpretation into an accepted canonical identifier.

Alias handling is part of local policy:

\begin{itemize}
\tightlist
\item
  an alias resolves to exactly one normalized reference at the relevant
  registry or policy version;
\item
  a missing alias, unsupported registry version, or ambiguous alias
  fails closed;
\item
  case folding, Unicode normalization, URI normalization, or prefix
  expansion is allowed only when the referenced identifier scheme
  defines it;
\item
  otherwise, normalized references are compared as exact strings;
\item
  set-like fields, such as scopes, resources, and authorization details,
  use the configured set mode after normalization.
\end{itemize}

Free-form natural-language intent, context, or capability text MAY be
carried as descriptive metadata. It MUST NOT be decision-authoritative
profile data.

\subsection{15. Production profile}\label{production-profile}

Production identity policy has two modes: disabled or required.

In disabled mode, the transport does not enforce L3 through L6 identity
policy. L0 through L2 lower-layer checks still apply according to the
configured TLS and attestation profile.

In required mode, the client fails closed when any required policy input
is missing, untrusted, expired, replayed, non-canonical, unsupported,
revoked, or inconsistent with the accepted TLS session. Peer-provided
metadata is never used as expected policy.

\begin{longtable}[]{@{}
  >{\raggedright\arraybackslash}p{(\columnwidth - 2\tabcolsep) * \real{0.5000}}
  >{\raggedright\arraybackslash}p{(\columnwidth - 2\tabcolsep) * \real{0.5000}}@{}}
\toprule\noalign{}
\begin{minipage}[b]{\linewidth}\raggedright
Area
\end{minipage} & \begin{minipage}[b]{\linewidth}\raggedright
Required-mode behavior
\end{minipage} \\
\midrule\noalign{}
\endhead
\bottomrule\noalign{}
\endlastfoot
Manager trust & Manager or policy-authority verification keys are
configured locally by \texttt{kid}, algorithm, issuer, audience, key
use, and public key, or loaded through an equivalent fail-closed key
source. \\
Algorithm safety & The token \texttt{alg} must match local key policy.
\texttt{none} is forbidden. Algorithm confusion is a hard failure. \\
Key separation & Manager signing keys and Agent confirmation keys are
separate trust domains and MUST NOT be accepted interchangeably. \\
Grant lifetime & Identity Grants are short-lived and carry \texttt{iat},
\texttt{exp}, and a unique \texttt{jti} or \texttt{cti}. \\
Grant revocation & Deployments that require early invalidation reject
revoked grant IDs and disabled Manager-key IDs before expiration. \\
Unknown keys & Unknown, disabled, retired, stale, wrong-use, or
mismatched key IDs fail closed. \\
Session-binding signer & Session Binding Statements are signed by the
confirmation key named in the verified grant, or by another key
explicitly authorized by that grant or local policy. \\
Session-binding fields & Session Binding Statements carry accepted
endpoint-key hash, request/exporter-context hash, one-shot nonce, and
attestation-binder hash when attestation binding is present. \\
Replay & JWT/CWT identity acceptance requires replay protection for
binding nonces or one-shot IDs. Multi-instance production deployments
use shared atomic insert-with-expiry semantics, such as
\texttt{SET\ NX\ EX}. \\
Local policy & Required L3 through L6 values are compared against local
expected policy. Missing expected values fail closed. \\
Error handling & Failure to check keys, revocation, session binding,
replay, canonicalization, or local identity policy is an authentication
failure, not a warning. \\
\end{longtable}

The initial mandatory JWT/JWS claim set is:

\begin{longtable}[]{@{}
  >{\raggedright\arraybackslash}p{(\columnwidth - 2\tabcolsep) * \real{0.5000}}
  >{\raggedright\arraybackslash}p{(\columnwidth - 2\tabcolsep) * \real{0.5000}}@{}}
\toprule\noalign{}
\begin{minipage}[b]{\linewidth}\raggedright
Token
\end{minipage} & \begin{minipage}[b]{\linewidth}\raggedright
Mandatory claims
\end{minipage} \\
\midrule\noalign{}
\endhead
\bottomrule\noalign{}
\endlastfoot
Identity Grant & \texttt{agtp\_type}, \texttt{agtp\_version},
\texttt{iss}, \texttt{aud}, \texttt{sub}, \texttt{jti}, \texttt{iat},
\texttt{exp}, \texttt{cnf.kid} \\
Session Binding Statement & \texttt{agtp\_type}, \texttt{agtp\_version},
\texttt{aud}, \texttt{jti}, \texttt{iat}, \texttt{exp},
\texttt{grant\_hash}, \texttt{leaf\_public\_key\_sha256},
\texttt{request\_context\_sha256}, \texttt{nonce}; also
\texttt{attestation\_binder\_sha256} when attestation binding is
present \\
\end{longtable}

When identity is accepted from JWT/JWS material, the Session Binding
Statement JWT is REQUIRED. The Identity Grant JWT authenticates
authority and semantics; it does not prove that the authorized Agent key
is present on the current TLS or exported-authenticator session. The
verifier accepts identity only after the Session Binding Statement JWT
is verified, linked to the exact signed Identity Grant JWT, compared
with the accepted session binding, checked for freshness and replay, and
compared with local expected policy.

CWT/COSE is a compact binary encoding for the same model. The grant
issuer must be trusted locally, the confirmation key must be named by
the verified grant, the COSE \texttt{kid} used for key lookup must be
protected, and the binding statement must bind the exact grant to the
accepted session.

\subsection{16. Freshness, replay, key rotation, and
revocation}\label{freshness-replay-key-rotation-and-revocation}

JWT/JWS and CWT/COSE adapters verify signatures against caller-provided
key lookup policy. They do not define how Manager keys are rotated, how
old keys are retired, how grants are revoked before expiration, or where
replay state is stored. Deployments own those decisions.

Deployments using this profile define:

\begin{itemize}
\tightlist
\item
  a trusted issuer namespace for Manager or policy-authority keys;
\item
  key identifiers, key-use rules, and key-version rules;
\item
  overlap windows for key rotation;
\item
  revocation sources for grants, Manager keys, or Agent binding keys
  when early invalidation is required;
\item
  maximum grant and session-binding lifetimes;
\item
  clock-skew tolerance;
\item
  replay-cache key shape;
\item
  replay-cache retention at least as long as the accepted
  session-binding lifetime.
\end{itemize}

Revocation has three separate targets:

\begin{itemize}
\tightlist
\item
  grant revocation rejects a specific Identity Grant by JWT \texttt{jti}
  or CWT \texttt{cti};
\item
  Manager-key revocation rejects grants signed by a disabled issuer key
  or \texttt{kid};
\item
  Agent binding-key revocation rejects Session Binding Statements signed
  by a compromised or retired confirmation key.
\end{itemize}

Unknown, revoked, stale, wrong-use, or ambiguous keys fail closed. A
deployment that relies on JWKS or another remote key set must define
freshness, cache lifetime, pinning or trust anchors, and failure
behavior for that key source.

Replay state is marked only after signature validation, session-binding
comparison, freshness checks, and local policy comparison succeed.
Failed attempts may be logged or rate-limited, but they do not consume a
one-shot nonce unless deployment policy deliberately chooses that
anti-probing behavior.

\subsection{17. Application response-cache
semantics}\label{application-response-cache-semantics}

Response caching and security-state caching are different things.

Response caching stores endpoint results for reuse. Replay caches store
one-shot freshness state. Verified Identity Grants, Session Binding
Statements, attestation evidence, authorization decisions, and prior
verification results MUST NOT be cached as acceptance evidence for a
later session or request.

Application response-cache controls, if used, follow RFC 9111's split
between shared and private caches, freshness, and response controls. The
default for profile-sensitive endpoints is \texttt{no-store}.

\begin{longtable}[]{@{}
  >{\raggedright\arraybackslash}p{(\columnwidth - 2\tabcolsep) * \real{0.5000}}
  >{\raggedright\arraybackslash}p{(\columnwidth - 2\tabcolsep) * \real{0.5000}}@{}}
\toprule\noalign{}
\begin{minipage}[b]{\linewidth}\raggedright
Directive
\end{minipage} & \begin{minipage}[b]{\linewidth}\raggedright
Meaning
\end{minipage} \\
\midrule\noalign{}
\endhead
\bottomrule\noalign{}
\endlastfoot
\texttt{no-store} & A cache MUST NOT store the response. This is
REQUIRED for identity tokens, Session Binding Statements, attestation
evidence, capability grants, authorization decisions, and responses
containing session-bound or caller-sensitive data. \\
\texttt{public} & A shared cache MAY store the response only when the
response is invariant across Agent ID, principal, tenant, authority
scope, policy grant, mTLS certificate identity, and session context. \\
\texttt{private} & A shared cache MUST NOT store the response. An
agent-local or principal-local cache MAY store it only inside the
relevant identity or policy partition. \\
\texttt{max-age} & Response freshness lifetime. It does not extend
token, grant, binding, revocation, attestation, or local policy
lifetimes. \\
\texttt{vary} & Identity, policy, or request fields that partition a
cached response when the response can differ by caller or authority
context. \\
\end{longtable}

\texttt{public} is a strong assertion. A public response must not depend
on caller identity, scope, principal, tenant, policy grant, mTLS
certificate, attestation result, session binding, task state, or any
other verifier-local policy input.

Private cache hits do not bypass current security checks. A verifier
still authenticates the current Identity Grant, verifies the current
Session Binding Statement, checks replay state, and compares local
policy before using a cached response. Cache hit is a performance
result, not an authorization result.

\subsection{18. Deployment topologies}\label{deployment-topologies}

\subsubsection{18.1 Direct-Agent mode}\label{direct-agent-mode}

In direct-Agent mode, the Agent terminates the TLS session and performs
the hardware-aware profile checks. The Agent signs the Session Binding
Statement with the confirmation key named by the verified Identity
Grant, or with an endpoint key explicitly authorized by the grant or
local policy.

\subsubsection{18.2 Gateway-routed mode}\label{gateway-routed-mode}

In gateway-routed mode, the gateway terminates the TLS session and
performs the hardware-aware profile checks. The gateway is the live TLS
endpoint. Gateway session binding proves the gateway endpoint, not the
final Agent process.

The deployment must also authenticate the gateway-to-Agent route before
treating the intended Agent as accepted. That can be modeled with
gateway ID, gateway public key, allowed route, and Agent route fields in
the Identity Grant, or with a separate gateway routing assertion.

A gateway routing assertion is not an Agent authority statement. It
authenticates the gateway's routing decision. The relying party still
needs an Identity Grant for the intended Agent and, when required, an
Agent-side holder-of-key proof for the gateway-to-Agent hop.

\subsection{19. Implementation hooks}\label{implementation-hooks}

The production client hook is:

\begin{itemize}
\tightlist
\item
  \texttt{atls.ClientConfig.IdentityPolicy} for local expected values;
\item
  \texttt{atls.ClientConfig.ObservedIdentity} for a session-bound
  observed assertion extracted from a trusted source.
\end{itemize}

The hook runs after exported-authenticator and attestation validation,
but before the accepted TLS connection is returned to the caller.

The client computes expected session binding from the accepted exported
authenticator: the leaf public key, the certificate request context, and
the attestation binder when attestation is present. The observed
assertion must carry matching binding values and a non-expired
\texttt{expires\_at} value before its identity fields are compared with
local policy.

The profile gives the client two defenses:

\begin{itemize}
\tightlist
\item
  relay defense: the observed assertion is tied to this accepted TLS
  session at L2;
\item
  diversion defense: the session-bound observed identity matches locally
  expected deployment, Agent, task, and authorization values at L3
  through L6.
\end{itemize}

\subsection{20. Validation algorithm}\label{validation-algorithm}

For each layer required by policy:

\begin{enumerate}
\def\labelenumi{\arabic{enumi}.}
\tightlist
\item
  Load the local expected value.
\item
  Reject if the expected value is missing, ambiguous, non-canonical, or
  only peer supplied.
\item
  Extract the observed session-bound assertion from an authenticated
  source.
\item
  Reject if the assertion is not bound to the accepted TLS session.
\item
  Reject if the assertion is expired or missing freshness metadata.
\item
  Extract the observed value for the required layer.
\item
  Reject if the observed value is missing or non-canonical.
\item
  Compare the observed value with the expected value.
\item
  Reject on mismatch.
\item
  Continue to the next required layer.
\end{enumerate}

For set-like values such as scopes, resources, or authorization details,
the observed set must satisfy local policy. The default
\texttt{contains\_all} mode accepts an observed set only when it
contains every locally required value. The stricter \texttt{exact} mode
also rejects extra observed values. A peer-provided scope list is not
enough by itself.

\subsection{21. Error handling and
diagnostics}\label{error-handling-and-diagnostics}

Validation failures preserve the layer, field, and error class. Callers
can distinguish local policy configuration errors from missing peer
claims, mismatched values, stale keys, replay, unsupported profile
versions, and binding mismatches.

Aggregated validation errors remain inspectable by layer and field so
callers can fail closed while still reporting actionable diagnostics.

The validator rejects unsafe identity-policy values, including invalid
UTF-8, control characters such as CRLF, and HTML delimiter characters.
It limits individual values to 1024 bytes and set-like fields to 128
values. Validation errors report only layer, field, and error class;
they do not echo raw peer values. HTTP or HTML presentation layers still
need normal output escaping and CSRF protections.

\subsection{22. Minimal implementation
path}\label{minimal-implementation-path}

A small implementation can be staged without changing the lower-layer
TLS or post-handshake attestation wire protocol:

\begin{enumerate}
\def\labelenumi{\arabic{enumi}.}
\tightlist
\item
  Define a local policy input structure for L3 through L6 expected
  values.
\item
  Define extraction points for observed values.
\item
  Authenticate the observed identity material before constructing an
  internal assertion.
\item
  Add fail-closed validators for exact-match string fields.
\item
  Add set-containment or exact-set validation for scopes, resources, and
  authorization details.
\item
  Add replay-cache support for one-shot binding nonces or IDs.
\item
  Wire validation before treating an accepted peer as the intended
  deployment, Agent, task, or authorized actor.
\end{enumerate}

The lower-layer verifier can remain focused on L0 through L2. L3 through
L6 enforcement belongs at the layer that has deployment policy, Agent
metadata, computation state, and authorization decisions.

\subsection{23. Implementation status}\label{implementation-status}

The reusable validator lives in \texttt{pkg/atls/identitypolicy}. It
implements the expected-versus-observed comparison model and
session-bound assertion validation described in this profile. The client
transport calls it when \texttt{atls.ClientConfig.IdentityPolicy} is
enabled.

The same package provides
\texttt{identitypolicy.NewAssertionFromSessionBinding} for the
post-authentication step. It builds an internal \texttt{Assertion} from
an already authenticated Identity Grant and an already-verified Session
Binding Statement. It does not parse wire tokens or verify signatures.
It checks that the statement is tied to the grant, audience, allowed
confirmation key, and minimum session-binding fields before policy
comparison. \texttt{identitypolicy.SessionBindingOptions} can attach a
replay cache for one-shot binding nonces.

The client config supports two integration styles. Callers can provide a
custom \texttt{ObservedIdentity} callback, or they can pass a verified
Identity Grant, a verified Session Binding Statement, and an optional
replay cache directly on \texttt{atls.ClientConfig}. In both cases, the
transport validates the resulting assertion against the accepted TLS
session before returning the connection. For the direct grant/statement
path, replay-cache marking happens only after session-binding and local
policy comparison succeed.

\texttt{pkg/agtp} provides concrete JWT/JWS and CWT/COSE wire-token
adapters. The JWT/JWS adapter verifies Identity Grants and Session
Binding Statements using locally configured issuer, audience, signing
methods, and key lookup policy. It does not choose trusted Manager keys,
rotate keys, perform revocation, define deployment policy, or own
distributed replay storage.

The initial JWT/JWS adapter treats grant \texttt{cnf.kid} as the
authorized session-binding signer key. The Session Binding Statement
signer is taken from the protected JWS \texttt{kid} header and later
checked by \texttt{identitypolicy} against the verified grant.
Thumbprint-based confirmation, such as \texttt{cnf.jkt}, can be added
when a deployment defines how key thumbprints map to local verification
keys.

For callers that want one fail-closed acceptance gate, \texttt{pkg/agtp}
exposes \texttt{VerifySessionIdentityJWT}. That helper verifies the
Manager-signed Identity Grant, verifies the Session Binding Statement,
checks binding-signer authorization, compares the resulting assertion
with local expected \texttt{identitypolicy.Policy} values and the
accepted TLS session binding, and only then marks the binding nonce in
the replay cache.

The implemented production profile covers:

\begin{itemize}
\tightlist
\item
  trusted Manager or policy-authority key lookup for Identity Grants
  through \texttt{JWTVerifyOptions};
\item
  issuer, audience, signing-method, \texttt{kid}, expiration, issued-at,
  token type, profile version, and \texttt{jti} checks;
\item
  confirmation-key binding through Identity Grant \texttt{cnf.kid};
\item
  Session Binding Statement signer authorization against the verified
  grant;
\item
  comparison with the accepted TLS endpoint key, request context, and
  optional attestation binder;
\item
  local \texttt{identitypolicy.Policy} comparison for required L3
  through L6 values;
\item
  replay-cache enforcement before the observed identity is accepted.
\end{itemize}

Deployment still chooses trusted keys, expected policy values,
revocation data, and distributed replay storage. Those sources can be
Manager configuration, Agent metadata, computation state, an
authorization policy engine, or a fail-closed registry integration. They
must not be raw peer-controlled metadata.

\subsection{Appendix A. Identity Grant JWT claim
map}\label{appendix-a.-identity-grant-jwt-claim-map}

The initial JWT/JWS profile uses two signed JWTs for different authority
questions:

\begin{itemize}
\tightlist
\item
  Identity Grant JWT: a Manager or policy authority signs upper-layer
  identity, task, and authorization values that the peer is allowed to
  claim.
\item
  Session Binding Statement JWT: the Agent confirmation key signs a
  holder-of-key statement binding the verified grant to the accepted TLS
  session.
\end{itemize}

The Session Binding Statement JWT does not authorize service, tenant,
task, scope, resource, or capability values. Those semantic values
belong in the Identity Grant JWT and are accepted only after local
policy comparison.

\begin{longtable}[]{@{}
  >{\raggedright\arraybackslash}p{(\columnwidth - 6\tabcolsep) * \real{0.2500}}
  >{\raggedright\arraybackslash}p{(\columnwidth - 6\tabcolsep) * \real{0.2500}}
  >{\raggedright\arraybackslash}p{(\columnwidth - 6\tabcolsep) * \real{0.2500}}
  >{\raggedright\arraybackslash}p{(\columnwidth - 6\tabcolsep) * \real{0.2500}}@{}}
\toprule\noalign{}
\begin{minipage}[b]{\linewidth}\raggedright
Identity Grant JWT field
\end{minipage} & \begin{minipage}[b]{\linewidth}\raggedright
Source or form
\end{minipage} & \begin{minipage}[b]{\linewidth}\raggedright
Purpose
\end{minipage} & \begin{minipage}[b]{\linewidth}\raggedright
Layer
\end{minipage} \\
\midrule\noalign{}
\endhead
\bottomrule\noalign{}
\endlastfoot
JWS protected header \texttt{kid} & Manager or policy-authority signing
key ID & Select the local verification key for the signed grant. The
value is only a key-selection hint. & prerequisite \\
JWS protected header \texttt{alg} & locally allowed signing method &
Prevent algorithm confusion. \texttt{none} is forbidden. &
prerequisite \\
\texttt{agtp\_type} & \texttt{agtp.identity-grant} & Distinguish the
grant token from other repository profile tokens. & profile guard \\
\texttt{agtp\_version} & \texttt{1} & Reject unsupported wire-profile
versions before interpreting claims. & profile guard \\
\texttt{iss} & trusted Manager or policy authority & Identify the
authority that issued the grant. & prerequisite \\
\texttt{aud} & relying service or application profile audience & Ensure
the grant was issued for this verifier or application profile. &
prerequisite \\
\texttt{sub} & upper-layer subject, usually the Agent when
\texttt{agent} is absent & Provide the subject authorized by the grant.
& L4 \\
\texttt{jti} & unique grant ID & Support grant uniqueness, revocation,
diagnostics, and replay-sensitive handling. & freshness \\
\texttt{iat} & issued-at time & Detect future-issued or malformed grants
and support freshness decisions. & freshness \\
\texttt{exp} & expiration time & Bound the lifetime of the authority
statement. & freshness \\
\texttt{cnf.kid} & Agent confirmation key ID & Name the key allowed to
sign the Session Binding Statement for this grant. & L4 / L2
prerequisite \\
\texttt{authorized\_endpoint\_keys} & optional endpoint key IDs & Allow
explicitly named endpoint keys when the deployment uses that binding
model. & L4 / L2 prerequisite \\
\texttt{service} & application service name or ID & Bind the grant to
the intended service. & L3 \\
\texttt{tenant} & tenant or account ID & Bind the grant to the intended
tenant. & L3 \\
\texttt{deployment} & deployment, region, cluster, or environment slice
& Bind the grant to the intended deployment target. & L3 \\
\texttt{environment} & production, staging, development, or similar &
Bind the grant to the intended environment class. & L3 \\
\texttt{workload} & workload, process, or component ID & Bind the grant
to the intended workload. & L4 \\
\texttt{agent} & Agent identity & Bind the grant to the intended Agent
rather than any peer on the same platform. & L4 \\
\texttt{agent\_public\_key} & stable Agent public-key reference or
fingerprint & Bind the Agent identity to a specific key when required by
policy. & L4 \\
\texttt{computation\_id} & computation or job ID & Bind the grant to a
computation instance. & L5 \\
\texttt{task\_id} & task ID & Bind the grant to the intended task. &
L5 \\
\texttt{thread\_id} & thread, conversation, or workflow ID & Bind the
grant to the intended thread or workflow context. & L5 \\
\texttt{delegation\_id} & delegation token or delegation record ID &
Bind the grant to a specific delegation chain or authorization handoff.
& L5 \\
\texttt{intent\_ref} & stable intent reference & Bind the grant to a
canonical task or intent family. & L5 \\
\texttt{capability\_ref} & stable capability reference & Bind the grant
to a canonical capability or resource-action class. & L6 \\
\texttt{ontology\_id} & stable policy vocabulary or namespace ID &
Identify the vocabulary used to interpret semantic references. & L6 \\
\texttt{scope} / \texttt{scopes} & OAuth-style scope values & State
coarse authorization strings for local policy comparison. & L6 \\
\texttt{resource} / \texttt{resources} & resource indicators or resource
URIs & State the resources the grant may be used against. & L6 \\
\texttt{authorization\_details} & profile-specific authorization detail
strings & Carry fine-grained authorization constraints. & L6 \\
\end{longtable}

\subsection{Appendix B. Session Binding Statement JWT claim
map}\label{appendix-b.-session-binding-statement-jwt-claim-map}

The Session Binding Statement JWT carries session facts, not semantic
authority. A verifier that accepts JWT/JWS identity material MUST verify
this JWT under the Agent confirmation key authorized by the verified
Identity Grant and MUST reject missing, mismatched, expired, replayed,
non-canonical, or unsupported binding material.

\begin{longtable}[]{@{}
  >{\raggedright\arraybackslash}p{(\columnwidth - 6\tabcolsep) * \real{0.2500}}
  >{\raggedright\arraybackslash}p{(\columnwidth - 6\tabcolsep) * \real{0.2500}}
  >{\raggedright\arraybackslash}p{(\columnwidth - 6\tabcolsep) * \real{0.2500}}
  >{\raggedright\arraybackslash}p{(\columnwidth - 6\tabcolsep) * \real{0.2500}}@{}}
\toprule\noalign{}
\begin{minipage}[b]{\linewidth}\raggedright
Session Binding Statement JWT field
\end{minipage} & \begin{minipage}[b]{\linewidth}\raggedright
Source or form
\end{minipage} & \begin{minipage}[b]{\linewidth}\raggedright
Purpose
\end{minipage} & \begin{minipage}[b]{\linewidth}\raggedright
Layer
\end{minipage} \\
\midrule\noalign{}
\endhead
\bottomrule\noalign{}
\endlastfoot
JWS protected header \texttt{kid} & Agent confirmation or explicitly
authorized endpoint key ID & Select the candidate binding-signature key.
The key is accepted only if authorized by the verified grant or local
policy. & L2/L4 prerequisite \\
JWS protected header \texttt{alg} & locally allowed signing method &
Prevent algorithm confusion. \texttt{none} is forbidden. &
prerequisite \\
\texttt{agtp\_type} & \texttt{agtp.session-binding} & Distinguish the
binding token from grants, envelopes, or unrelated JWTs. & profile
guard \\
\texttt{agtp\_version} & \texttt{1} & Reject unsupported wire-profile
versions before interpreting binding fields. & profile guard \\
\texttt{aud} & relying service or application profile audience & Ensure
the statement was issued for this verifier or profile. & L2
prerequisite \\
\texttt{jti} & unique binding statement ID & Support diagnostics and
one-shot replay handling. & freshness \\
\texttt{iat} & issued-at time & Detect future-issued or malformed
binding statements. & freshness \\
\texttt{exp} & expiration time & Bound the lifetime of the session
proof. It should not exceed the grant lifetime. & freshness \\
\texttt{grant\_hash} & domain-separated hash of exact signed Identity
Grant bytes & Bind the statement to the exact Manager-signed grant that
was verified. & L2 \\
\texttt{leaf\_public\_key\_sha256} & hash of accepted TLS or
exported-authenticator endpoint public key & Bind the statement to the
accepted endpoint key. & L2 \\
\texttt{request\_context\_sha256} & hash of accepted
request/exported-authenticator context & Bind the statement to the
accepted request context. & L2 \\
\texttt{attestation\_binder\_sha256} & hash of accepted attestation
binder, when present & Bind the statement to accepted
attestation-to-channel evidence. & L2 \\
\texttt{nonce} & one-shot nonce or unique binding value & Prevent replay
of an otherwise valid binding statement. & freshness \\
\end{longtable}

The receiver compares these fields as strict, already-canonical values.
It does not repair aliases, reserialize JSON to compute
\texttt{grant\_hash}, infer missing binding fields from peer metadata,
or treat a valid Identity Grant JWT as session-bound without the Session
Binding Statement JWT.

\subsection{Appendix C. Authority separation and envelope
packaging}\label{appendix-c.-authority-separation-and-envelope-packaging}

A deployment may package the Manager-signed Identity Grant inside, or
alongside, the Agent-signed Session Binding Statement so an
implementation passes around one envelope. Packaging does not merge
authorities. The verifier still validates the Manager or
policy-authority signature on the grant, validates the Agent
confirmation-key signature on the session binding, checks that the
binding names the exact verified grant, and compares the resulting
assertion with local policy.

Moving semantic grant fields directly into an Agent-signed Session
Binding Statement would make the Agent the authority for service,
tenant, task, scope, resource, and capability claims. That changes the
trust model and is unsafe for profiles that need Manager- or
policy-authority-issued authorization.

A single Manager-signed JWT that also includes session-binding values is
possible only when the Manager participates in, or is given trustworthy
evidence of, the per-connection TLS binding. That design adds an online
authority dependency and moves replay and freshness responsibilities to
the Manager path.

The two-JWT design adds one additional JWS verification plus a
domain-separated hash over the exact compact Identity Grant JWT bytes.
This is usually small compared with TLS setup and attestation
verification. Verifiers MUST NOT cache a verified Identity Grant as
acceptance evidence across sessions. Internal optimization must still
re-check signature, expiration, revocation state, profile version,
issuer, audience, grant hash, fresh Session Binding Statement, replay
state, and local policy for each accepted session.

The reusable \texttt{pkg/agtp} JWT/JWS adapter also supports a
single-envelope verification path. In that profile, the outer envelope
is signed by the Agent binding key and carries the exact inner
Manager-signed Identity Grant JWT plus the session-binding fields. The
verifier still authenticates the inner Manager grant, authenticates the
outer Agent binding statement, compares the outer \texttt{grant\_hash}
with the exact inner grant bytes, enforces replay and freshness, and
compares the resulting assertion with local policy. Agent-signed
semantic claims in the outer envelope are not Manager authorization. The
current runtime client configuration remains wired to the two-token
JWT/JWS profile unless a caller explicitly uses the envelope verifier.

\subsection{Appendix D. Non-normative trust-flow
example}\label{appendix-d.-non-normative-trust-flow-example}

\begin{Shaded}
\begin{Highlighting}[]
\NormalTok{sequenceDiagram}
\NormalTok{  participant M as Bob Manager}
\NormalTok{  participant B as Bob Agent}
\NormalTok{  participant A as Alice}
\NormalTok{  participant X as Malware Agent on Alice machine}

\NormalTok{  M{-}\textgreater{}\textgreater{}B: Identity Grant signed by Manager key}
\NormalTok{  Note over M,B: Grant authorizes Bob Agent key}

\NormalTok{  A{-}\textgreater{}\textgreater{}B: Start TLS or hardware{-}aware TLS session}
\NormalTok{  B{-}\textgreater{}\textgreater{}A: Identity Grant}
\NormalTok{  B{-}\textgreater{}\textgreater{}A: Session Binding Statement signed by Bob Agent key}

\NormalTok{  X{-}{-}\textgreater{}\textgreater{}A: Injects peer{-}provided identity metadata}
\NormalTok{  X{-}{-}\textgreater{}\textgreater{}A: Replays another grant or binding statement}

\NormalTok{  A{-}\textgreater{}\textgreater{}A: Verify grant with trusted Bob Manager public key}
\NormalTok{  A{-}\textgreater{}\textgreater{}A: Verify session binding with Bob Agent key from grant}
\NormalTok{  A{-}\textgreater{}\textgreater{}A: Check expected service, Agent, task, and capability}
\NormalTok{  A{-}\textgreater{}\textgreater{}A: Reject metadata not signed or not bound to this session}
\end{Highlighting}
\end{Shaded}

Bob Manager issues an authorization grant for Bob Agent. Bob Agent
communicates with Alice and binds the Manager grant to the current TLS
or hardware-aware TLS session. Malware on Alice's machine can try to
inject peer-provided metadata or reuse an old grant or binding
statement, but it cannot create a valid Manager-signed grant or
Bob-Agent-key-signed session binding for the current accepted session.

Alice accepts only the Manager-signed grant and the Agent-signed session
binding after local expected service, Agent, task, and capability checks
pass. This does not protect a fully compromised Alice process or local
policy store.

\end{document}
